% TEMPLATE-MILC-v3.tex - plantilla maestra UNICA de las guias MILC v3.
%
% La usan las cuatro familias de guias, cada una con su driver:
%   · Grados 6-11          -> scripts/build-guias-g11.py
%   · Bebras               -> scripts/build-guias-bebras.py
%   · Semillero            -> scripts/build-guias-semillero.py
%   · Territorio interior  -> scripts/build-guias-territorio-interior.py
%
% Antes habia tres copias casi identicas (~400 lineas iguales, 6 distintas):
% cada mejora habia que aplicarla tres veces, y en la practica no se hacia
% (el motor de recursos vivia solo en dos de ellas). Lo que varia entre
% familias esta parametrizado en 6 placeholders que cada driver llena
% (se nombran aqui SIN su sintaxis real: el builder reemplaza por texto plano,
% tambien dentro de los comentarios, y un valor multilinea rompe el preambulo):
%
%   HEADER_IZQ        encabezado izquierdo de pagina
%   HEADER_DER        encabezado derecho de pagina
%   PORTADA_ETIQUETA  rotulo sobre el numero grande (GRADO/MOMENTO/MODULO)
%   PORTADA_SERIE     linea de serie de la portada
%   PORTADA_ID        identificador de la guia en la portada
%   PORTADA_CIRCULO   el circulito de grado (°), o vacio si no aplica
%   PDF_TITULO        titulo en texto plano (sin \textbf ni comillas TeX) para
%                     los metadatos del PDF (/Title); lo produce lib_xelatex.texto_pdf
%
% Composicion (auditoria 2026-09-03): las bandas de estacion piden espacio
% (\Needspace*) y prohiben el salto justo despues (\nopagebreak), asi no quedan
% huerfanas al pie; las cajas largas (softbox, drawbox, infoband, puentebox) son
% `breakable`, asi una caja de media pagina ya no arrastra todo a la siguiente
% dejando la anterior al 40 %. Todo texto sobre mostaza o verde va en milcNegro
% (blanco sobre #E5B400 da 1,93:1 y sobre #58A923 2,95:1; ninguno pasa WCAG AA).
% El resto de claves las documenta content/guias/_SCHEMA.md. El script valida
% que no queden placeholders sin reemplazar antes de escribir el archivo final.

% Etiquetado PDF (latex-lab): /Lang, estructura y texto alternativo de las
% figuras, para lectores de pantalla. Debe ir ANTES de \documentclass.
\DocumentMetadata{lang=es-CO,tagging=on}
\documentclass[11pt,letterpaper]{article}
% headheight=24pt: la cabecera derecha lleva dos lineas (identidad de la guia
% y estacion actual). headsep se acorta en la misma medida para que el bloque
% de texto quede exactamente donde estaba con headheight=12pt.
\usepackage[margin=2.0cm,headheight=24pt,headsep=13pt]{geometry}
\usepackage[table]{xcolor}
\usepackage{fontspec}
\usepackage[spanish]{babel}
\usepackage{tikz}
% Librerías TikZ para los diagramas de las guías (bloque `recursos:`):
%  · babel  → babel en español vuelve activos caracteres como «>», y TikZ los usa
%             en sus opciones (>=stealth, ->). Sin esta librería, cualquier
%             diagrama con flechas revienta con «\language@active@arg> has an
%             extra }». Debe ir ANTES de que se dibuje nada.
%  · positioning        → sintaxis «below left=2pt and 3pt».
%  · decorations.pathreplacing → llaves ({brace}) para señalar rangos.
\usetikzlibrary{babel,positioning,decorations.pathreplacing}
\usepackage{graphicx}
% Circuitos: el trabajo de electrónica se hace hoy con micro:bit y sus sensores
% integrados (MakeCode, sin cableado), así que no se necesitan esquemáticos.
% Si algún día entran montajes con protoboard, basta descomentar esta línea:
% los diagramas .tex/.tikz declarados en `recursos:` ya se incrustan solos.
% \usepackage{circuitikz}
\usepackage[most]{tcolorbox}
\usepackage{tabularx}
\usepackage{array}
\usepackage{fancyhdr}
\usepackage{titlesec}
\usepackage[document]{ragged2e}  % bandera a la derecha en todo el cuerpo
\usepackage{enumitem}
\usepackage{microtype}
\usepackage{etoolbox}
\usepackage{needspace}
% hyperref al final del preambulo: metadatos del PDF (titulo, autor, idioma,
% familia), marcadores por estacion (\pdfbookmark) y enlaces sin recuadro.
\usepackage[hidelinks,
  pdftitle={Edición de imagen — lo que se puede tocar y lo que no},
  pdfauthor={PhD. Álvaro Cárdenas Orozco},
  pdfsubject={Tecnología e Informática · Grado 8 · Guía 24 · MILC v3},
  pdflang={es-CO},
  bookmarksopen=true,bookmarksnumbered=false]{hyperref}

% Helvetica Neue no trae la flecha U+2192, asi que cada «→» escrito literalmente
% en el YAML se compilaba a NADA, en silencio (462 flechas en 61 guias: rutas del
% tipo «paleta → tipografia → lockup» salian rotas). Mapeamos el caracter a su
% version matematica, que si tiene glifo. Asi el autor puede seguir escribiendo
% «→» comodamente en el YAML.
\usepackage{newunicodechar}
\newunicodechar{→}{\ensuremath{\rightarrow}}
% Mismo problema con los simbolos de marca y vinetas: el «✓» de «una palomita ✓»
% salia como cuadrito vacio en el PDF de todas las guias que lo usaban.
\usepackage{pifont}
\newunicodechar{✓}{\ding{51}}
\newunicodechar{✗}{\ding{55}}
\newunicodechar{✔}{\ding{52}}
\newunicodechar{•}{\textbullet}
\newunicodechar{≥}{\ensuremath{\geq}}
\newunicodechar{≤}{\ensuremath{\leq}}

\setmainfont{Helvetica Neue}
\setsansfont{Helvetica Neue}
\setlength{\parindent}{0pt}
\setlength{\parskip}{6pt}
% Sin particion de palabras: en 6.º-7.º una palabra cortada por guion al final
% de linea («Comuni-cacion») frena la lectura mucho mas que una linea corta.
\hyphenpenalty=10000
\exhyphenpenalty=10000
\pretolerance=2000
\tolerance=3000
\setlength{\emergencystretch}{3em}
\linespread{1.10}
\renewcommand{\arraystretch}{1.25}
\setlist{leftmargin=5mm,itemsep=1mm,topsep=1mm}
\newcolumntype{Y}{>{\RaggedRight\arraybackslash}X}

\definecolor{milcMagenta}{HTML}{E6007E}
\definecolor{milcVino}{HTML}{5A0038}
\definecolor{milcTurquesa}{HTML}{0093A5}
\definecolor{milcVerde}{HTML}{58A923}
\definecolor{milcMostaza}{HTML}{E5B400}
\definecolor{milcOcre}{HTML}{B86600}
\definecolor{milcNegro}{HTML}{241816}
\definecolor{milcGris}{HTML}{F7F7F4}
\definecolor{milcLinea}{HTML}{E8E2DD}
\definecolor{milcGradePrimary}{HTML}{FF6600}
\definecolor{milcGradeDark}{HTML}{9A3E00}
\definecolor{milcGradeSoft}{HTML}{FFE4D0}
\definecolor{milcGradeAccent}{HTML}{CFE7FF}
\definecolor{milcGradeText}{HTML}{FFFFFF}
\color{milcNegro}

\pagestyle{fancy}
\fancyhf{}
% Cabecera de dos lineas: identidad de la guia arriba y, a la derecha y en
% gris, la estacion en curso (\markright desde \milcestacion). Antes de la
% primera estacion la segunda linea va vacia; el \strut mantiene la altura
% para que la linea de arriba no baile entre paginas.
\lhead{\small Tecnología e Informática · Grado 8\\[-1pt]{\footnotesize\strut}}
\rhead{\small Guía 24 · MILC v3\\[-1pt]{\footnotesize\strut\color{milcNegro!65}\rightmark}}
\cfoot{\small\thepage}
\renewcommand{\headrulewidth}{0pt}

% Sobre mostaza (#E5B400) y verde (#58A923) el texto va en milcNegro; sobre el
% resto de la paleta, en blanco. Se usa dentro de fonttitle/fontupper, donde un
% \color posterior manda sobre coltitle.
\newcommand{\milctintasobre}[1]{%
  \ifboolexpr{test{\ifstrequal{#1}{milcMostaza}} or test{\ifstrequal{#1}{milcVerde}}}%
    {\color{milcNegro}}{\color{white}}}

\newtcolorbox{titlebox}[1]{
  enhanced,colback=#1,colframe=#1,arc=7mm,boxrule=0pt,
  left=7mm,right=7mm,top=3mm,bottom=3mm,
  before skip=8pt,after skip=8pt,
  fontupper=\sffamily\bfseries\Large\color{white}
}

\newtcolorbox{softbox}[3]{
  enhanced,breakable,pad at break=3mm,
  colback=#2,colframe=#1,borderline west={4pt}{0pt}{#1},
  arc=2mm,boxrule=.4pt,left=4mm,right=4mm,top=3mm,bottom=3mm,
  before skip=6pt,after skip=8pt,title={#3},coltitle=white,
  colbacktitle=#1,fonttitle=\sffamily\bfseries\milctintasobre{#1},fontupper=\RaggedRight,
  attach boxed title to top left={xshift=3mm,yshift=-2mm},
  boxed title style={arc=2mm, boxrule=0pt}
}

\newtcolorbox{drawbox}[2]{
  enhanced,breakable,pad at break=4mm,
  colback=white,colframe=#1,arc=4mm,boxrule=1pt,
  left=5mm,right=5mm,top=4mm,bottom=4mm,before skip=8pt,after skip=8pt,
  title={#2},coltitle=white,colbacktitle=#1,fonttitle=\sffamily\bfseries\milctintasobre{#1},
  fontupper=\RaggedRight,
  attach boxed title to top left={xshift=4mm,yshift=-2mm},
  boxed title style={arc=3mm, boxrule=0pt}
}

\newtcolorbox{infoband}[1]{
  enhanced,breakable,pad at break=2mm,colback=#1!8,colframe=#1!50,arc=2mm,boxrule=.5pt,
  left=4mm,right=4mm,top=2mm,bottom=2mm,before skip=4pt,after skip=4pt,
  fontupper=\small\RaggedRight
}

% Puente narrativo entre secciones (contrato editorial Conéctate).
% Da continuidad: "Acabas de X. Ahora vas a Y porque Z."
% Visualmente: banda izquierda en color del grado, texto en itálica pequeña.
\newtcolorbox{puentebox}{
  enhanced,breakable,pad at break=2mm,colback=milcGradeSoft,colframe=milcGradeDark,
  borderline west={3pt}{0pt}{milcGradeDark},
  arc=0mm,boxrule=0pt,left=5mm,right=4mm,top=2.5mm,bottom=2.5mm,
  before skip=4pt,after skip=8pt,
  fontupper=\itshape\small\color{milcNegro}\RaggedRight
}
% La flecha va en modo matematico a proposito: Helvetica Neue NO tiene el glifo
% U+2192, asi que \textrightarrow se compilaba a NADA («Missing character: There
% is no → in font Helvetica Neue») y el puente salia sin flecha. En modo
% matematico la flecha viene de la fuente matematica y si se dibuja.
\newcommand{\puente}[1]{%
  \begin{puentebox}%
    \textcolor{milcGradeDark}{$\rightarrow$}\hspace{2mm}#1%
  \end{puentebox}%
}

\newcommand{\linead}[1]{\noindent\color{milcLinea}\rule{#1}{0.5pt}\color{milcNegro}}
\newcommand{\respuesta}[1]{\par\vspace{2mm}\foreach \n in {1,...,#1}{\noindent\color{milcLinea}\rule{\linewidth}{0.5pt}\par\vspace{5mm}}\color{milcNegro}}
\newcommand{\checkbox}{\textcolor{milcGradeDark}{\fbox{\phantom{x}}}\hspace{5pt}}

% ─── Listas aireadas para las guías socioemocionales ────────────────────
% Componer las verificaciones como «\checkbox texto\\» deja el texto que salta
% de línea debajo del cuadrito y apelmaza el bloque. Estos entornos alinean la
% continuación y airean los ítems. Los usa Territorio Interior; cualquier
% familia puede hacerlo.
\newlist{checklista}{itemize}{1}
\setlist[checklista]{
  label=\textcolor{milcGradeDark}{\fbox{\phantom{x}}},
  leftmargin=9mm, labelsep=3.2mm, itemsep=2.4mm,
  topsep=2.5mm, parsep=0pt, partopsep=0pt, before=\RaggedRight
}
\newlist{pasoslista}{enumerate}{1}
\setlist[pasoslista]{
  label=\textcolor{milcGradeDark}{\sffamily\bfseries\arabic*.},
  leftmargin=8mm, labelsep=3mm, itemsep=2.8mm,
  topsep=1mm, parsep=0pt, partopsep=0pt, before=\RaggedRight
}

\newcommand{\phasecard}[4]{
\begin{tcolorbox}[enhanced,colback=#3,colframe=#2,arc=3mm,boxrule=.5pt,height=3.05cm,valign=center,halign=center,left=1.5mm,right=1.5mm,top=2mm,bottom=2mm]
{\sffamily\bfseries\fontsize{22}{24}\selectfont\textcolor{#2}{#1}}\par
{\sffamily\bfseries\small #4}
\end{tcolorbox}
}

% ── Piezas del rediseno 2026 (legibilidad 6.º-11.º) ───────────────────
% Banda de estacion: reemplaza el titlebox macizo. Titulo a la izquierda,
% dato practico (tiempo/modalidad) a la derecha, en una sola linea.
% Nunca huerfana: pide 9 lineas libres (o salta de pagina rellenando con
% \vfill) y prohibe el salto justo despues de la banda. Nueve y no seis:
% la caja partible que sigue necesita ~80pt (titulo adosado, relleno y dos
% lineas) para arrancar en la misma pagina; con menos, tcolorbox la manda
% entera a la siguiente y la banda se queda sola. El penalty 10000 va
% en `after`, ANTES del espacio posterior: un `after skip` (aunque sea 0pt)
% es un \vskip tras la caja, y ese glue es un punto de corte legal que un
% \nopagebreak escrito despues de \end{tcolorbox} ya no cierra. Cada banda
% deja marcador PDF y actualiza la cabecera derecha.
\newcounter{milcestacion}
\newcommand{\milcestacion}[2]{%
\Needspace*{9\baselineskip}%
\stepcounter{milcestacion}%
\pdfbookmark[1]{#2}{milcestacion\themilcestacion}%
\markright{#2}%
\begin{tcolorbox}[enhanced,colback=#1,colframe=#1,arc=2mm,boxrule=0pt,
  left=5mm,right=5mm,top=2.6mm,bottom=2.6mm,before skip=11pt,
  after={\par\nopagebreak\vskip7pt}]
{\sffamily\bfseries\fontsize{15}{18}\selectfont\milctintasobre{#1}#2}
\end{tcolorbox}}

% Tarjeta de paso numerada: sustituye las tablas «Pilar / Aplicacion», que
% eran formato de consulta docente, no de estudiante. El numero va en un
% circulo sobre el borde para que la secuencia se lea de un vistazo.
\newcommand{\milcpaso}[3]{%
\Needspace{4\baselineskip}%
\begin{tcolorbox}[enhanced,colback=white,colframe=milcLinea,arc=1.5mm,boxrule=.9pt,
  left=14mm,right=4mm,top=3mm,bottom=3mm,before skip=4pt,after skip=4pt,
  fontupper=\RaggedRight,
  overlay={\node[anchor=north west,circle,fill=#2,text=white,inner sep=0pt,
    minimum size=7.4mm,font=\sffamily\bfseries\small]
    at ([xshift=3.6mm,yshift=-3.2mm]frame.north west) {#1};}]
#3
\end{tcolorbox}}

% Casilla marcable de tamano util con lapiz (la anterior era \fbox{\phantom{x}},
% de alto variable segun el texto que la seguia).
\newcommand{\milcmarca}{\framebox[3.8mm]{\rule{0pt}{3.4mm}}}

% Fila marcable de lista suelta (checks de la estacion 1).
\newcommand{\milccheckrow}[1]{%
\par\vspace{1.8mm}\noindent
\makebox[6.4mm][l]{\raisebox{-.55mm}{\textcolor{milcNegro}{\milcmarca}}}%
\parbox[t]{\dimexpr\linewidth-6.4mm\relax}{\RaggedRight #1}\par}

% Celda marcable dentro de la rubrica (tabla de autoevaluacion). Misma
% sangria francesa, pero pensada para vivir dentro de una columna X.
\newcommand{\milcopcion}[1]{%
\makebox[5.2mm][l]{\raisebox{-.55mm}{\textcolor{milcNegro}{\milcmarca}}}%
\parbox[t]{\dimexpr\linewidth-5.2mm\relax}{\RaggedRight #1}}

% ── Recursos visuales de la guía ──────────────────────────────────────
% Declarados en el bloque `recursos:` del YAML y emitidos por el builder.
% \guiaFigura   → imagen rasterizada o PDF (foto, carta celeste, montaje).
% \guiaDiagrama → fuente TikZ/circuitikz (.tex) incrustada como vector:
%                 es la vía para circuitos y esquemáticos editables.
% [#1] = texto alternativo (alt) · #2 = ruta relativa al .tex · #3 = pie
% El alt llega al PDF etiquetado (/Alt) y lo lee el lector de pantalla.
\newcommand{\guiaFigura}[3][]{%
\begin{center}
\includegraphics[width=0.88\linewidth,height=0.38\textheight,keepaspectratio,alt={#1}]{#2}\\[1.5mm]
{\footnotesize\itshape\textcolor{milcNegro!75}{#3}}
\end{center}
}

\newcommand{\guiaDiagrama}[2]{%
\begin{center}
\input{#1}\\[1.5mm]
{\footnotesize\itshape\textcolor{milcNegro!75}{#2}}
\end{center}
}

\begin{document}
\thispagestyle{empty}

% ── Portada ───────────────────────────────────────────────────────────
% Antes: un rectangulo de color a pagina completa. Se veia bien en pantalla,
% pero en la fotocopiadora del colegio se comia el toner y salia gris sucio.
% Ahora la banda ocupa el tercio superior y el resto va en blanco.
\begin{tikzpicture}[remember picture,overlay]
  \path[fill=milcGradePrimary]
    (current page.north west) rectangle ([yshift=-10.6cm]current page.north east);

  \node[anchor=north west,text=milcGradeText] at ([xshift=2.0cm,yshift=-1.55cm]current page.north west)
    {\sffamily\bfseries\fontsize{12}{14}\selectfont GRADO};
  \node[anchor=north west,text=milcGradeText,opacity=.85] at ([xshift=2.0cm,yshift=-2.35cm]current page.north west)
    {\sffamily\bfseries\fontsize{8.5}{10}\selectfont SERIE GUÍAS · TECNOLOGÍA E INFORMÁTICA};
  \node[anchor=north east,text=milcGradeText,opacity=.85,align=right] at ([xshift=-2.0cm,yshift=-1.55cm]current page.north east)
    {\sffamily\bfseries\fontsize{9}{12}\selectfont I.E. Sor María Juliana\\Cartago, Valle del Cauca};

  \node[anchor=west,text=milcGradeText] (gradonum) at ([xshift=1.85cm,yshift=-7.1cm]current page.north west)
    {\sffamily\bfseries\fontsize{112}{112}\selectfont 8};
  % El circulito de grado (°) solo aplica a las guias de grado («11°»); en
  % Bebras y Semillero el numero grande es un momento/modulo, no un grado.
  % Se posiciona relativo al nodo `gradonum`, no en coordenadas absolutas.
  \draw[milcGradeText,line width=1.6pt]
    ([xshift=2.0mm,yshift=-4.5mm]gradonum.north east) circle (.15cm);

  \node[anchor=west,text=milcGradeText,text width=11.4cm,align=left] at ([xshift=1.5cm,yshift=0.5cm]gradonum.east)
    {
     {\sffamily\bfseries\fontsize{9}{11}\selectfont\MakeUppercase{Período 3 · Multimedia, narrativa y ciberseguridad}}\\[3mm]
     {\sffamily\bfseries\fontsize{21}{25}\selectfont\setlength{\baselineskip}{27pt}Lo que se puede tocar\\[4pt]y lo que no}\\[3.5mm]
     {\sffamily\bfseries\fontsize{15}{18}\selectfont Guía 24-8-TIC}
    };
\end{tikzpicture}

\vspace*{9.4cm}
\vfill

{\sffamily\bfseries\fontsize{8}{10}\selectfont\textcolor{milcGradeDark}{LO QUE TE LLEVAS HOY}}

\begin{tcolorbox}[enhanced,colback=white,colframe=milcNegro,arc=2mm,boxrule=1.6pt,
  left=5mm,right=5mm,top=4mm,bottom=4mm,before skip=3pt,after skip=12pt,
  fontupper=\RaggedRight\fontsize{11}{15.5}\selectfont]
Una imagen real editada con cuatro ajustes básicos, recorte, brillo, contraste y texto, en una herramienta gratuita, con la captura del antes y del después. Y la autoevaluación ética en cinco preguntas: si hay un rostro identificable, si tienes permiso, si la edición cambia el sentido, si declaras la fuente y si la imagen final trata bien a quien aparece.
\end{tcolorbox}

\begin{tcolorbox}[enhanced,colback=milcGris,colframe=milcGris,arc=2mm,boxrule=0pt,
  left=5mm,right=5mm,top=3.5mm,bottom=3.5mm,after skip=12pt]
{\sffamily\bfseries\fontsize{8}{10}\selectfont\textcolor{milcNegro!55}{ESTA GUÍA ES DE}}\\[3mm]
\begin{tabularx}{\linewidth}{X p{3.2cm} p{3.2cm}}
\linead{\linewidth} & \linead{3.0cm} & \linead{3.0cm}\\[-1mm]
{\fontsize{8.5}{10}\selectfont\textcolor{milcNegro!55}{Nombre completo}} &
{\fontsize{8.5}{10}\selectfont\textcolor{milcNegro!55}{Grupo}} &
{\fontsize{8.5}{10}\selectfont\textcolor{milcNegro!55}{Fecha}}\\
\end{tabularx}
\end{tcolorbox}

\vfill

\noindent\textcolor{milcLinea}{\rule{\linewidth}{1pt}}\\[2mm]
{\sffamily\bfseries\fontsize{9.5}{12}\selectfont Autor: Dr. Álvaro Cárdenas Orozco}\\[1mm]
{\sffamily\fontsize{8.5}{11}\selectfont\textcolor{milcNegro!55}{Metodología MILC · apertura ancestral · pensamiento computacional · triángulo Dussel–estoicismo–Floridi}}

\newpage

\section*{Apertura: Edición de imagen --- lo que se puede tocar y lo que no}

\begin{softbox}{milcGradeDark}{milcGradeSoft}{Saber ancestral}
Hace más de mil años, orfebres del Cauca medio, hoy Quindío, fundieron en oro poporos y figuras que los coleccionistas llamaron «quimbayas». En 1891 el Estado colombiano compró 122 de esas piezas, encontradas en Filandia. En 1893 el presidente encargado Carlos Holguín las regaló a la reina regente de España. Nadie les preguntó a los pueblos del Cauca medio ni a los colombianos. Las piezas siguen en el Museo de América, en Madrid. En 2017 la Corte Constitucional dijo que ese regalo violó dos derechos de todos: la moralidad administrativa y el patrimonio público. Y ordenó al Estado hacer todo lo necesario para «repatriar y/o readquirir las 122 piezas» (Corte Constitucional, 2017). Lo que está en juego: algo que no era de quien lo regaló, sacado del lugar que le daba sentido. La \textbf{cara de exclusión}: ocho años después, las piezas no han vuelto. La Cancillería alegó que ninguna ley internacional obliga a devolver un regalo entre Estados. Y los orfebres que las hicieron no tienen voz en el pleito. Hoy vas a editar una imagen y antes de tocarla vas a preguntarte lo mismo: de quién es y quién dio permiso.\par\smallskip{\footnotesize\textcolor{milcNegro!70}{\textbf{Fuente.} Corte Constitucional de Colombia. (2017). \emph{Sentencia SU-649/17} (M. P. Alberto Rojas Ríos), 19 de octubre.}}
\end{softbox}

\begin{softbox}{milcTurquesa}{milcGris}{Saber contemporáneo}
Editar una imagen es sencillo: Canva, Photopea, GIMP o la app del celular hacen los \textbf{cuatro ajustes básicos}. \textbf{Recortar} quita partes para enfocar lo importante. \textbf{Brillo} aclara u oscurece. \textbf{Contraste} separa lo claro de lo oscuro para que se vea mejor. \textbf{Texto} agrega un título o una nota. Con eso se resuelve casi todo. Lo difícil no es la herramienta: es la \textbf{ética visual}. La Ley 1581 de 2012 protege los datos personales, y la imagen del rostro es uno de ellos: usarla exige permiso de la persona, o de sus padres si es menor. Y la edición puede mentir sin cambiar un píxel del rostro: un recorte que quita el contexto o un texto que la persona nunca dijo cambian el sentido. Por eso, antes de publicar, cinco preguntas: ¿hay rostro identificable?, ¿tengo permiso?, ¿la edición cambia el sentido?, ¿declaro la fuente?, ¿la imagen final trata bien a quien aparece? Si una responde mal, no se publica.
\end{softbox}

\begin{softbox}{milcMostaza}{milcGris}{Pregunta-puente}
\textit{Las piezas de Filandia se regalaron sin preguntarle a nadie y se exhiben lejos de donde tenían sentido. Si recortas una foto de un compañero y le pones un texto que no dijo, ¿qué le estás quitando y a quién le pediste permiso?}
\end{softbox}

\begin{softbox}{milcVerde}{milcGris}{Saber y saber hacer}
Al terminar podrás: (1) \textbf{identificar} cómo un recorte o un texto cambian el sentido de la misma foto; (2) \textbf{aplicar} los cuatro ajustes básicos en una herramienta gratuita; (3) \textbf{evaluar} tu edición con las cinco preguntas éticas; (4) \textbf{explicar} por qué el permiso y la fuente no son opcionales.
\end{softbox}



\small
\begin{tabularx}{\textwidth}{>{\bfseries}p{3.0cm}Y}
\rowcolor{milcGradePrimary!12}Autor & Dr. Álvaro Cárdenas Orozco\\
DBA & Producción multimedia ética y comportamiento responsable en entornos digitales (MEN, T\&I)\\
Referentes & Dussel (estética de la liberación) · Ley 1336 · Floridi (infoesfera)\\
Producto final & Una imagen real editada con cuatro ajustes básicos, recorte, brillo, contraste y texto, en una herramienta gratuita, con la captura del antes y del después. Y la autoevaluación ética en cinco preguntas: si hay un rostro identificable, si tienes permiso, si la edición cambia el sentido, si declaras la fuente y si la imagen final trata bien a quien aparece.\\
\end{tabularx}
\normalsize

\puente{En la sesión 3 hablaste con voz propia. Hoy tocas la imagen de alguien, y la regla cambia: la voz es tuya, la imagen es de quien aparece. En la sesión 5 vas a trabajar con audio.}

\milcestacion{milcGradeDark}{Ruta MILC: así vamos a aprender}
\begin{tabularx}{\textwidth}{>{\centering\arraybackslash}X>{\centering\arraybackslash}X>{\centering\arraybackslash}X>{\centering\arraybackslash}X}
\phasecard{1}{milcVerde}{milcGris}{Escucha\\[-1mm]\small Observo, escucho y pregunto.} &
\phasecard{2}{milcTurquesa}{milcGris}{Entender\\[-1mm]\small Sistematizo y ordeno.} &
\phasecard{3}{milcMagenta}{milcGris}{Hacer\\[-1mm]\small Construyo y aplico.} &
\phasecard{4}{milcVino}{milcGris}{Cierre\\[-1mm]\small Evalúo y reflexiono.}
\end{tabularx}


\puente{Antes de editar, vas a comparar tres versiones de la misma foto y ver cómo un recorte o un texto le cambian el sentido sin tocar el rostro.}

\milcestacion{milcVerde}{Estación 1 de 3 · Escucha}

\begin{softbox}{milcVerde}{milcGris}{Escena para escuchar}
\textbf{Actividad 1 · IDENTIFICA --- Tres versiones de la misma foto} (15 min · individual). (1) Tu docente proyecta tres versiones de una misma foto: la original, una recortada que deja por fuera algo importante, y una con un texto agregado que cambia el sentido. (2) Para cada versión escribe en una línea qué crees que está pasando en la foto. (3) Compara tus tres líneas: ¿cuál versión te hizo pensar algo distinto de lo que muestra la original? (4) Escribe qué se quitó en el recorte y qué se agregó en el texto. (5) Escribe en una línea a quién habría que pedirle permiso antes de publicar cualquiera de las tres.
\end{softbox}

{\sffamily\bfseries\fontsize{8}{10}\selectfont\textcolor{milcNegro!55}{MARCA CUANDO LO HAYAS HECHO}}
\milccheckrow{Escribí qué pasa en cada una de las tres versiones.}
\milccheckrow{Identifiqué qué quitó el recorte y qué agregó el texto.}
\milccheckrow{Escribí a quién habría que pedirle permiso.}

\begin{infoband}{milcVerde}


\textbf{Lo que va al cuaderno (Actividad 1 · IDENTIFICA).} \textbf{Título:} «Actividad 1 --- Tres versiones de la misma foto». \textbf{Formato:} tres líneas, una por versión, la nota de qué se quitó y qué se agregó, y la línea del permiso. \textbf{Extensión:} media página. \textbf{Sabes que terminaste cuando} puedes decir cuál versión miente y con qué herramienta lo hace, recorte o texto.
\end{infoband}

\puente{Ya viste cómo se miente con un recorte. Ahora aprendes los cuatro ajustes y las cinco preguntas, y con tu pareja editas una imagen con permiso.}

\milcestacion{milcTurquesa}{Estación 2 de 3 · Entender}

\begin{softbox}{milcTurquesa}{milcGris}{Lo que vas a entender}
\textbf{Actividad 2 · APLICA --- Los cuatro ajustes con permiso} (30 min · parejas). (1) Con tu pareja, elijan una foto de las que trae tu docente, sin rostros de personas reales. O una que tomó uno de los dos, donde solo aparece esa persona y da permiso en voz alta. (2) Abran la foto en Canva, Photopea o la app del celular y guarden una copia del original: ese es el antes. (3) Recorten para enfocar lo importante sin quitar el contexto que cambia el sentido. (4) Ajusten brillo y contraste hasta que se vea mejor sin que parezca otra cosa. (5) Agreguen un texto corto que describa, no que invente, y guarden el después.
\end{softbox}

\milcpaso{1}{milcTurquesa}{\textbf{Los cuatro ajustes.} Recorte: quita lo que sobra, pero si quita el contexto, miente. Brillo: aclara u oscurece toda la imagen. Contraste: separa claros y oscuros; poco deja todo gris, mucho quema los detalles. Texto: describe lo que hay; si pone palabras en boca de alguien, miente.}
\milcpaso{2}{milcTurquesa}{\textbf{Las cinco preguntas.} ¿Hay rostro identificable? Si sí, hace falta permiso explícito. ¿Tengo permiso? Oral o escrito, de la persona o de sus padres si es menor. ¿La edición cambia el sentido? Si un recorte o un texto hacen que parezca otra cosa, no. ¿Declaro la fuente? Si la foto no es tuya, se dice de quién es. ¿La imagen final trata bien a quien aparece? Si esa persona la viera, ¿se sentiría engañada?}
\milcpaso{3}{milcTurquesa}{\textbf{La lección de las 122 piezas.} Sacadas del Cauca medio y exhibidas en Madrid, siguen siendo del lugar que les dio sentido. Una foto recortada y publicada sin permiso hace lo mismo con una persona: la saca de su contexto y la exhibe como si fuera tuya.}
\milcpaso{4}{milcTurquesa}{\textbf{Cómo se hace, paso a paso.} (1) Elegir una foto con permiso. (2) Guardar el antes. (3) Recortar sin quitar contexto. (4) Brillo y contraste sin exagerar. (5) Texto que describe. (6) Guardar el después. (7) Cinco preguntas.}

\begin{drawbox}{milcTurquesa}{La edición con ética, en una mirada}
\textbf{Recorte:} enfoca, pero no quita contexto. \\[1mm] \textbf{Brillo y contraste:} se ve mejor, no parece otra cosa. \\[1mm] \textbf{Texto:} describe, no inventa. \\[1mm] \textbf{Permiso:} de quien aparece, o de sus padres. \\[1mm] \textbf{Fuente:} de quién es la foto.
\end{drawbox}

\begin{softbox}{milcMostaza}{milcGris}{Errores comunes y su corrección}
(1) \textbf{Editar sin permiso}: usar la foto de un compañero sin pedírsela. (2) \textbf{Recortar para engañar}: dejar por fuera lo que cambiaría la lectura. (3) \textbf{Texto que pone palabras}: una frase que la persona no dijo. (4) \textbf{Foto ajena sin fuente}: bajarla de internet y no decir de quién es. (5) \textbf{Filtro que cambia el cuerpo}: una edición que la persona no reconocería como suya.
\end{softbox}

\begin{infoband}{milcTurquesa}


\textbf{Lo que va al cuaderno (Actividad 2 · APLICA).} \textbf{Título:} «Actividad 2 --- Los cuatro ajustes con permiso». \textbf{Formato:} el nombre de la foto y de quién es, los cuatro ajustes aplicados con una línea cada uno sobre qué cambió, y la nota del permiso. \textbf{Extensión:} media página. \textbf{Sabes que terminaste cuando} tienes el antes y el después guardados y puedes decir qué hizo cada ajuste.
\end{infoband}

\puente{Ya sabes editar. Toca elegir tu propia imagen, aplicar los cuatro ajustes, guardar el antes y el después y responder las cinco preguntas por escrito.}

\milcestacion{milcMagenta}{Estación 3 de 3 · Hacer}

\begin{softbox}{milcMagenta}{milcGris}{Cómo vas a construir tu producto}
\textbf{Actividad 3 · EVALÚA --- Tu imagen y las cinco preguntas} (25 min · individual). (1) Elige tu propia imagen: una que tomaste tú sin rostros ajenos, o una con permiso dado en voz alta y anotado. (2) Aplica los cuatro ajustes y guarda el antes y el después. (3) Responde por escrito las cinco preguntas, con una línea de explicación cada una. (4) Si alguna respuesta es «no» o «no sé», escribe qué tendrías que hacer antes de publicar. (5) Muéstrale el después a la persona que aparece, o a tu pareja si no aparece nadie. Pregúntale si reconoce la foto y si la publicaría.
\end{softbox}

\milcpaso{1}{milcMagenta}{\textbf{Tu entrega tiene tres piezas.} (1) El antes y el después, exportados o con enlace. (2) Las cinco preguntas respondidas con su línea de explicación. (3) Lo que dijo la persona que aparece, o tu pareja, al ver el después.}
\milcpaso{2}{milcMagenta}{\textbf{Una respuesta mala detiene todo.} Si no hay permiso, no se publica. Si el recorte cambia el sentido, se vuelve al original. Si la fuente no se sabe, se busca o no se usa. Las cinco preguntas no son un promedio: una sola basta para parar.}
\milcpaso{3}{milcMagenta}{\textbf{El permiso se anota.} «Me dijo que sí» se olvida. «Le pregunté a Sara el martes y dijo que sí, para el afiche del festival» es un permiso. Para menores, el permiso lo dan los padres.}
\milcpaso{4}{milcMagenta}{\textbf{La prueba de quien aparece.} La persona ve tu edición. Si dice «esa soy yo, publícala», pasaste. Si dice «yo no dije eso» o «ahí no se ve que estábamos jugando», el recorte o el texto mintieron.}

\begin{drawbox}{milcMagenta}{Antes de entregar}
\checkbox Antes y después guardados y exportados. \\[1mm] \checkbox Cuatro ajustes aplicados: recorte, brillo, contraste y texto. \\[1mm] \checkbox El recorte no quita contexto ni el texto inventa. \\[1mm] \checkbox Permiso anotado con nombre y fecha, o foto sin rostros ajenos. \\[1mm] \checkbox Fuente declarada si la foto no es mía. \\[1mm] \checkbox Cinco preguntas respondidas con una línea cada una. \\[1mm] \checkbox La persona que aparece, o mi pareja, vio el después y dijo si la publicaría.
\end{drawbox}

\begin{softbox}{milcMagenta}{milcGris}{Plantilla de trabajo}
\textbf{Foto.} \textit{De …, tomada el …} \\[1mm] \textbf{Permiso.} \textit{Le pregunté a … el … y dijo …} \\[1mm] \textbf{Recorte.} \textit{Quité … y dejé … porque …} \\[1mm] \textbf{Brillo y contraste.} \textit{Subí / bajé … para …} \\[1mm] \textbf{Texto.} \textit{Escribí «…», que describe …} \\[1mm] \textbf{Quien aparece dijo.} \textit{…}
\end{softbox}

\begin{infoband}{milcMagenta}


\textbf{Lo que va al cuaderno (Actividad 3 · EVALÚA).} \textbf{Título:} «Actividad 3 --- Tu imagen y las cinco preguntas». \textbf{Formato:} el antes y el después pegados o con enlace, las cinco preguntas con su respuesta y su línea, la nota del permiso y lo que dijo quien aparece. \textbf{Extensión:} una página. \textbf{Sabes que terminaste cuando} las cinco respuestas están escritas y ninguna es «no» sin un plan para arreglarlo.
\end{infoband}

\puente{Lo que entregas hoy: el antes y el después, y la autoevaluación de cinco preguntas. La prueba de calidad: la persona que aparece ve tu edición y no se siente engañada.}

\milcestacion{milcMostaza}{Producto de la sesión}
\textbf{Imagen editada con cuatro ajustes + antes y después + cinco preguntas éticas respondidas + permiso anotado}.
\begin{softbox}{milcMostaza}{milcGris}{Criterios de calidad}
(1) El antes y el después están guardados y se ven los cuatro ajustes. (2) El recorte no quita contexto que cambie el sentido y el texto describe, no inventa. (3) Hay permiso anotado con nombre y fecha, o la foto no tiene rostros ajenos. (4) La fuente está declarada si la foto no es tuya. (5) Las cinco preguntas tienen respuesta y una línea de explicación cada una. (6) Quien aparece, o tu pareja, vio el después y dijo si la publicaría.
\end{softbox}

\puente{Antes de cerrar, mira lo que hiciste desde cinco lados. Una edición modesta con permiso vale más que una espectacular que alguien no autorizó.}

\milcestacion{milcVino}{Evaluación liberadora · cinco dimensiones}

\begin{softbox}{milcVino}{milcGris}{Sentido de la evaluación}
Evaluar no es solo revisar una nota. Es reconocer qué aprendí, qué cambió en mí, qué emociones manejé, cómo me relaciono con la ciudad y la comunidad, y qué le dejaré a quien viene detrás.
\end{softbox}

\textbf{Mi reflexión liberadora en cinco dimensiones.} Completa cada frase iniciada:

\small
\begin{tabularx}{\textwidth}{>{\bfseries\RaggedRight\arraybackslash}p{3.4cm}Y>{\RaggedRight\arraybackslash}p{4.6cm}}
\rowcolor{milcVino}\color{white}Dimensión & \color{white}\bfseries Mi respuesta (frame starter) & \color{white}\bfseries Algo que descubrí\\
1. Personal & Lo que más cambió en mí fue \linead{6.5cm} & \linead{4.2cm}\\
2. Emocional & La emoción más fuerte fue \linead{2.5cm} y la regulé \linead{3cm} & \linead{4.2cm}\\
3. Ciudadana & En democracia esto importa porque \linead{6cm} & \linead{4.2cm}\\
4. Local (Cartago) & En mi barrio sería útil para \linead{6.5cm} & \linead{4.2cm}\\
5. Intergeneracional & A mi abuela/abuelo le diría \linead{6.5cm} & \linead{4.2cm}\\
\end{tabularx}
\normalsize

\begin{infoband}{milcVino}
\textbf{Para la próxima sustentación.} Vuelve a esta tabla en seis meses. Verás que las respuestas cambian — no porque lo aprendido cambie, sino porque tú cambias. La evaluación liberadora es una conversación contigo a través del tiempo.
\end{infoband}


\puente{Para terminar, tres ideas sobre la imagen de otro, contadas con palabras de esta guía: Dussel sobre el rostro que no es instrumento, Epicteto sobre juzgar sin saber, y Floridi sobre lo que somos hechos de información.}

\milcestacion{milcGradeDark}{Tres ideas para cerrar · Dussel, un estoico y Floridi}

\begin{softbox}{milcVino}{milcGris}{Enrique Dussel · lente del nosotros}
\textit{El rostro del otro se revela como alguien, no como una cosa más entre mis instrumentos.}

{\footnotesize\itshape\color{milcNegro!70}--- idea tomada de Enrique Dussel, contada con palabras de esta guía. No es una cita textual.}

Una foto de tu compañero en tu carpeta se siente como un archivo más. No lo es: es un rostro, alguien que puede decir «yo no dije eso». Pedir permiso es tratar la imagen como rostro y no como instrumento.

\textbf{¿Qué foto de alguien tengo guardada como si fuera un archivo cualquiera?}
\end{softbox}
\linead{\linewidth}\\[1mm]
\linead{\linewidth}

\begin{softbox}{milcMostaza}{milcGris}{Estoicismo · Epicteto · cuidado interior}
\textit{Si ves a alguien hacer algo, no digas que lo hizo mal: no sabes desde dónde juzgar.}

{\footnotesize\itshape\color{milcNegro!70}--- idea tomada de Epicteto, contada con palabras de esta guía. No es una cita textual.}

Un recorte deja por fuera que estaban jugando; un texto pone una frase que nadie dijo. Quien mira la foto editada juzga sin saber. La ética visual empieza por no fabricar ese juicio.

\textbf{¿Qué foto recortada me hizo juzgar a alguien esta semana, y qué no me mostró?}
\end{softbox}
\linead{\linewidth}\\[1mm]
\linead{\linewidth}

\begin{softbox}{milcTurquesa}{milcGris}{Luciano Floridi · lente de la infoesfera}
\textit{Somos organismos hechos también de información, que vivimos en un entorno de información.}

{\footnotesize\itshape\color{milcNegro!70}--- idea tomada de Luciano Floridi, contada con palabras de esta guía. No es una cita textual.}

Tu imagen es parte de ti, no un adorno que se puede tomar. Editar la foto de alguien sin permiso es tocar a esa persona. Las piezas de Filandia en Madrid muestran lo que pasa cuando algo así se hace a escala de un país.

\textbf{¿Qué parte de mí anda en fotos que otros editaron sin preguntarme?}
\end{softbox}
\linead{\linewidth}\\[1mm]
\linead{\linewidth}

\begin{infoband}{milcNegro}
\textbf{Nota para el docente.} Las tres ideas están contadas con palabras de esta guía a partir del pensamiento de cada autor; no son citas textuales. De 9.º en adelante se trabajan con la obra y su referencia.
\end{infoband}

\milcestacion{milcVino}{Mi autoevaluación · marca una casilla por fila}

{\small
\setlength{\extrarowheight}{2.2pt}
\begin{tabularx}{\textwidth}{>{\RaggedRight\arraybackslash\bfseries}p{3.0cm}YYY}
\rowcolor{milcVino}\color{white}\bfseries Criterio & \color{white}\bfseries Lo logré & \color{white}\bfseries En proceso & \color{white}\bfseries Necesito apoyo\\[.6mm]
Comprensión del tema & \milcopcion{Explico las ideas con mis palabras.} & \milcopcion{Reconozco las ideas, falta precisión.} & \milcopcion{Necesito repasar lo básico.}\\[.8mm]
\rowcolor{milcGris}Pensamiento computacional & \milcopcion{Usé los pilares al armar mi producto.} & \milcopcion{Usé algunos pilares.} & \milcopcion{Necesito releer la estación 2.}\\[.8mm]
Aplicación y producto & \milcopcion{Mi producto cumple los criterios.} & \milcopcion{Está armado, con ajustes pendientes.} & \milcopcion{Aún está en construcción.}\\[.8mm]
\rowcolor{milcGris}Reflexión 5 dimensiones & \milcopcion{Respondí cada una con honestidad.} & \milcopcion{Respondí algunas con detalle.} & \milcopcion{Mis respuestas son muy generales.}\\[.8mm]
Triángulo de pensamiento & \milcopcion{Conecté las 3 voces con mi trabajo.} & \milcopcion{Conecté al menos una.} & \milcopcion{No las relaciono con mi proceso.}\\[.8mm]
\end{tabularx}}

\vspace{3mm}
\textbf{Hoy entendí que:}\\[1mm]
\linead{\linewidth}\\[1mm]
\linead{\linewidth}

\textbf{Mi compromiso para la próxima sesión es:}\\[1mm]
\linead{\linewidth}\\[1mm]
\linead{\linewidth}

\begin{softbox}{milcMostaza}{milcGris}{Cita de despedida · Marco Aurelio}
\textit{``El obstáculo en el camino se vuelve el camino. Lo que estorba al camino se vuelve camino.''}\\[1mm]
Cada error de hoy, bien observado, se convierte en pieza del camino que pisarás mañana.
\end{softbox}

\small
\begin{tabularx}{\textwidth}{>{\bfseries}p{3.6cm}Y}
\rowcolor{milcGradeDark!12}Nombre completo: & \linead{10cm}\\
Fecha: & \linead{4cm} \quad Curso/grupo: \linead{4cm}\\
Mi firma: & \linead{9cm}\\
\end{tabularx}
\normalsize

\begin{infoband}{milcGradeDark}
\textbf{Para la próxima sesión.} Llega con tu antes, tu después y las cinco preguntas. En la sesión 5 vas a grabar y editar audio: recorte, volumen y ruido, y a exportar en dos formatos. La regla del permiso sigue: la voz de alguien también es suya.
\end{infoband}

\end{document}
